% =====================================================================
%  K-LD :: The semantics, its fidelity (RQ1), and mechanised lemmas (RQ4).
%  The differential study (RQ2/RQ3) is in _sec_e3.tex.
% =====================================================================
\providecommand{\kld}{\textsc{K-LD}}
\providecommand{\kst}{\textsc{K-ST}}
\providecommand{\cmark}{\ensuremath{\checkmark}}

\section{The \texorpdfstring{\kld{}}{K-LD} Semantics and Its Validation}
\label{sec:semantics}

We present \kld{} and answer two questions here: is it a \emph{faithful} model of
\gls{iec}~61131-3 \gls{ld} (RQ1, \S\ref{sec:eval:rq1}), and can its
per-construct correctness be \emph{machine-checked} (RQ4, \S\ref{sec:eval:rq4})?
The remaining questions---whether \kld{} works as an oracle and how sensitively it
catches translation defects (RQ2/RQ3)---are answered by the differential study of
\S\ref{sec:e3}.

\subsection{The Semantics in \texorpdfstring{\textsf{K}}{K}}
\label{sub:sem}
\kld{} is a configuration of cells rewritten scan by scan. The retained state is a
single Boolean \emph{image} ($\langle\code{image}\rangle$) mapping every input,
coil, and function-block done-bit to its current value; alongside it,
$\langle\code{timers}\rangle$ and $\langle\code{counters}\rangle$ hold per-instance
function-block state, and $\langle\code{dt}\rangle$ is the cycle time~$\Delta t$.
The driver cell $\langle\code{k}\rangle$ runs the scan loop over an input trace
$\langle\code{inputs}\rangle$ (one Boolean assignment per scan), logging the image
after each scan to $\langle\code{trace}\rangle$. One scan overlays the next input,
solves the program's rungs top to bottom against the image, and appends the result:
\[
\text{\code{\#run}} \Rightarrow P \,\rightsquigarrow\, \text{\code{\#endScan}}
\,\rightsquigarrow\, \text{\code{\#run}}
\qquad\text{(consuming one input, per scan).}
\]
Contacts read the image ($\code{XIC}(X)\Rightarrow \code{image}[X]$;
$\code{XIO}(X)\Rightarrow\lnot\,\code{image}[X]$), series and parallel compose by
conjunction and disjunction, and \code{OTE}/\code{OTL}/\code{OTU} write, set, or
retain their coil. Because only recomputed coils change each scan while latches and
function-block state persist in their cells, cross-scan retention---the crux of
timing and sequencing---is modelled directly. The timers are a faithful port of the
standard state machine: each carries $(\mathrm{ET},\mathrm{STATE},\mathrm{prevIN})$
with $\mathrm{STATE}\in\{\text{idle},\text{timing},\text{done}\}$, and elapsed time
is measured from the trigger edge ($\mathrm{ET}=0$ on the edge scan), so a
\code{TON} fires exactly $\lceil \mathrm{PT}/\Delta t\rceil$ scans after its input
rises. From this one definition \textsf{K} yields both the interpreter (\code{krun})
used below and the prover (\code{kprove}) used in \S\ref{sec:eval:rq4}.

\subsection{RQ1: An Executable Semantics Faithful to IEC~61131-3}
\label{sec:eval:rq1}

\paragraph{Method: validation against OpenPLC.} The behavioural reference for the
function blocks is the C that OpenPLC actually executes: \gls{matiec}'s reference
implementations (\code{iec\_std\_FB.h}) of \code{TON}, \code{TOF}, \code{TP},
\code{CTU}, and \code{CTD}. We extract those bodies, compile a harness that drives
each block with a controlled clock (elapsed time advanced by a fixed $\Delta t$ per
scan) and a chosen input trace, and record the per-scan done-bit. We then run \kld{}
on the same program and input trace and compare the two traces scan for scan. This
executes the \emph{same} reference code deployed on open-hardware \gls{plc}s, so
agreement means \kld{} matches real \gls{plc} behaviour, not merely a paper
specification.

\paragraph{Result: exact agreement.} Table~\ref{tab:kld:conformance} shows the
outcome for the five function blocks (preset~$3$ for timers, $2$ for counters,
$\Delta t=1$). \kld{} reproduces the OpenPLC/\gls{matiec} done-bit trace
\emph{exactly, scan for scan}, in every case. The combinational and latch
constructs are likewise confirmed: an \code{AND} rung follows its contacts in both
polarities, and a set/reset latch yields the retention trace
$\langle\top,\top,\bot,\bot\rangle$---the light stays latched after its set input
drops, exactly as \code{OTL}/\code{OTU} require.

\begin{table}[t]
  \centering
  \caption{RQ1---\kld{} vs the OpenPLC/\gls{matiec} reference, done-bit trace per
    scan (timers $\mathrm{PT}=3$, counters $\mathrm{PV}=2$, $\Delta t=1$; $\top$/$\bot$
    for true/false). \kld{} matches the executed reference on every scan of every
    block.}
  \label{tab:kld:conformance}
  \small
  \begin{tabular}{@{}llc@{}}
    \toprule
    Block & Done-bit trace (\kld{} $=$ OpenPLC) & Match \\
    \midrule
    \code{TON} (on-delay)  & $\bot\,\bot\,\bot\,\top\,\bot$                 & \cmark \\
    \code{TOF} (off-delay) & $\top\,\top\,\top\,\top\,\top\,\bot$           & \cmark \\
    \code{TP}  (pulse)     & $\bot\,\top\,\top\,\top\,\bot\,\top$           & \cmark \\
    \code{CTU} (count up)  & $\bot\,\bot\,\bot\,\bot\,\top\,\bot\,\bot$     & \cmark \\
    \code{CTD} (count down)& $\bot\,\bot\,\bot\,\bot\,\top\,\top\,\bot$     & \cmark \\
    \bottomrule
  \end{tabular}
\end{table}

\paragraph{The validation is discriminating.} That the check is meaningful---not a
tautology---is shown by the defects it caught while \kld{} was being written. An
early version of all three timers fired \emph{one scan too early}: it counted the
trigger scan itself as one $\Delta t$ of elapsed time, whereas the standard (and
\gls{matiec}) measure elapsed time from the trigger edge. A second version let
\code{CTU} increment past its preset instead of saturating. Both diverged from the
reference and were flagged by exactly this comparison; the timers were then
rewritten as a direct port of the reference state machine, and \code{CTU} was
capped, after which the traces coincide as in Table~\ref{tab:kld:conformance}. A
validation that catches such off-by-one and saturation errors is a genuine test of
fidelity, and it is what licenses using \kld{} as the reference oracle in
\S\ref{sec:e3}.

\subsection{RQ4: Machine-Checked Construct Lemmas}
\label{sec:eval:rq4}

RQ1 grounds \kld{} by \emph{execution}; RQ4 goes further and \emph{proves}, in
\code{kprove}, that \kld{}'s rules implement the \gls{iec}~61131-3 input/output
relation for each construct. We emphasise scope: we do not mechanise a full
\kld{}$\leftrightarrow$GOTO equivalence, which would require formalising
\gls{esbmc}'s \gls{ir} as well; the translation is validated empirically in
\S\ref{sec:e3}, and RQ4 instead furnishes machine-checked correctness of the
\emph{oracle itself}.

\paragraph{Per-construct correctness.} For a construct we state a reachability
claim: from an image with symbolic bit values, executing the construct's rung
yields exactly the standard-mandated coil value. Each is discharged against the
semantics by \code{kprove} (Haskell backend). Table~\ref{tab:kld:proof} lists the
seven lemmas proved for the combinational and latch fragment---contacts (series,
parallel, negation), \code{OTE}, and \code{OTL}/\code{OTU} in both their acting and
their \emph{retaining} cases---all discharged (\code{kprove} returns
$\top$).

\begin{lemma}[Per-construct correctness, representative case]
\label{lem:kld:construct}
For all identifiers and all bit values $v_a,v_b$, executing the rung
$\code{OTE}(q)\!:=\!\code{XIC}(a)\ \code{*}\ \code{XIC}(b)$ from an image binding
$a\!\mapsto\!v_a,\ b\!\mapsto\!v_b$ terminates with $q\mapsto v_a\wedge v_b$; the
analogous statements hold for parallel contacts ($\vee$), \code{XIO} ($\lnot$), and
for \code{OTL}/\code{OTU} setting/resetting on a true rung and \emph{retaining} the
coil on a false one.
\end{lemma}

\begin{table}[t]
  \centering
  \caption{RQ4---construct-correctness lemmas machine-checked in \code{kprove}. All
    seven discharge ($\top$).}
  \label{tab:kld:proof}
  \small
  \begin{tabular}{@{}llc@{}}
    \toprule
    Lemma & Property proved & \code{kprove} \\
    \midrule
    L1 & series contacts $\Rightarrow$ coil $= \code{bit}(a)\wedge\code{bit}(b)$ & \cmark \\
    L2 & parallel contacts $\Rightarrow$ coil $= \code{bit}(a)\vee\code{bit}(b)$  & \cmark \\
    L3 & \code{XIO} $\Rightarrow$ coil $= \lnot\,\code{bit}(a)$                    & \cmark \\
    L4 & \code{OTL} sets the coil on a true rung                                  & \cmark \\
    L5 & \code{OTL} \emph{retains} the coil on a false rung                       & \cmark \\
    L6 & \code{OTU} resets the coil on a true rung                                & \cmark \\
    L7 & \code{OTU} \emph{retains} the coil on a false rung                       & \cmark \\
    \bottomrule
  \end{tabular}
\end{table}

\paragraph{Boundary of the mechanised guarantee.} The timer and counter constructs
are not yet mechanised: their correctness is a multi-scan property requiring
induction over the input list, and the current scan-loop encoding is not discharged
automatically by \code{kprove} (the prover does not narrow the symbolic input list
into the constructor forms the scan rule matches, so the induction does not start).
We state this plainly as the boundary of the proof: the combinational and latch
fragment is \emph{machine-checked}, while the timer and counter constructs are
\emph{execution-validated} against OpenPLC (RQ1) and \emph{empirically validated}
through the differential study (\S\ref{sec:e3}). Closing the inductive lemmas---by
a list case-split lemma or a narrowing-friendly reformulation of the scan
driver---is the natural next step, and would lift the full fragment to a mechanised
guarantee.
